% Setup Table S2b -- held-fixed training/evaluation hyperparameters. \input{} from report-draft.tex Ch6.
\begin{table}[tbp]
\centering
\small
\begin{tabular}{l l l}
\toprule
\textbf{Hyperparameter} & \textbf{Default (unless stated)} & \textbf{Varied / swept in} \\
\midrule
Alignment weight $\lambda$ & 0.5 (after Yu et al.~\cite{REPA}) & Appendix~\ref{ch:appendix} \\
Similarity                 & cosine                            & --- \\
Token averaging            & per residue                       & \S\ref{sec:repa} (design knob) \\
Residue length $n$         & 256 (headline regime)             & 128 (\S\ref{sec:proteina-scale}) \\
Training regime            & PDB \emph{and} AFDB (both reported) & --- \\
Sampler                    & SDE, noise scale $\gamma = 0.45$  & \S\ref{sec:proteina-sampler} \\
Evaluation                 & supercolumn suite (Table~\ref{tab:setup-metrics}) & --- \\
\bottomrule
\end{tabular}
\caption{\textbf{Held-fixed training and evaluation hyperparameters: our ``unless stated otherwise'' defaults.} The alignment weight follows the value Yu et al.~\cite{REPA} found effective
for image diffusion; cosine similarity averaged per residue is the REPA objective of
\S\ref{sec:repa}, with the per-residue (vs.\ per-sample) averaging mode being a design
knob we discuss there. We anchor headline claims at residue length $n{=}256$ and use
$n{=}128$ as a scale-robustness check (\S\ref{sec:proteina-scale}). Every variant is
trained in both data regimes and sampled at the default SDE noise scale $\gamma{=}0.45$,
which we ablate across six values in \S\ref{sec:proteina-sampler}. Optimisation
hyperparameters (batch size, learning rate, step counts) and the $\lambda$/batch-size/projector-depth
ablations are in Appendix~\ref{ch:appendix}.}
\label{tab:setup-hparams}
\end{table}
